\chapter{Implicit Differentiation and Related Rates}

So far almost every function we have differentiated arrived in a friendly form: a rule that takes an
input $x$, does something to it, and hands back a single output $y = f(x)$. You feed in a number,
you turn the crank, out comes the answer. The derivative measures how the crank's output responds
when you nudge the input.

But the world is not always so polite. Some of the most important curves in geometry and physics
refuse to be written as $y = f(x)$. A circle is the cleanest example: the equation
\[
x^2 + y^2 = 25
\]
describes a perfectly definite curve, and that curve obviously has tangent lines everywhere, with
perfectly definite slopes. Yet there is no single function $f$ with $y = f(x)$ whose graph is the
whole circle, because most vertical lines hit the circle \emph{twice}. The circle is a relationship
between $x$ and $y$, not a machine that turns $x$ into $y$.

This chapter is about two closely linked ideas. The first is \emph{implicit differentiation}: a way
to find slopes on curves that are tied together by an equation, without ever solving for $y$. We will
see that it is not a new rule at all -- it is the chain rule, worn as a disguise. The second is
\emph{related rates}: a method for problems where two or more changing quantities are locked together
by a constraint, so that the rate of one drives the rate of another. A ladder slides, a balloon
inflates, a shadow lengthens; in each case a single equation ties the pieces together, and
differentiating that equation in time tells us how the rates relate. Both ideas spring from the same
root: \emph{when quantities are bound by an equation, their rates of change are bound too.}

\section{Curves That Are Not Functions}
\label{ch13:not-functions}

Let us look hard at the circle $x^2 + y^2 = 25$. Pick a point on it, say $(3,4)$. Geometrically the
circle has an unambiguous tangent line there. We can even guess its slope: a radius runs from the
center $(0,0)$ out to $(3,4)$ with slope $4/3$, and the tangent is perpendicular to the radius, so
the tangent slope should be $-3/4$. The geometry leaves no doubt that a slope exists. The only
question is how to \emph{compute} it from the equation.

\begin{figure}[H]
\centering
\begin{tikzpicture}
\begin{axis}[
  width=10cm, height=10cm,
  axis equal image,
  axis lines=middle,
  xlabel={$x$}, ylabel={$y$},
  xmin=-6.5, xmax=6.5, ymin=-6.5, ymax=6.5,
  xtick={-5,5}, ytick={-5,5},
  samples=120,
]
\addplot[curveblue, thick, domain=0:360] ({5*cos(x)}, {5*sin(x)});
% point (3,4)
\addplot[only marks, mark=*, tangentred] coordinates {(3,4)};
% tangent line at (3,4): slope -3/4
\addplot[tangentred, thick, domain=-1:6.5] {4 - (3/4)*(x-3)};
% radius
\addplot[structurepurple, dashed, thick] coordinates {(0,0) (3,4)};
\node[tangentred] at (axis cs: 5.3,1.2) {tangent};
\node[structurepurple] at (axis cs: 0.9,2.6) {radius};
\node at (axis cs: 3.55,4.4) {$(3,4)$};
\end{axis}
\end{tikzpicture}
\caption{The circle $x^2+y^2=25$ is not the graph of any single function of $x$, yet it has a
definite tangent line at each point. At $(3,4)$ the radius has slope $4/3$, so the perpendicular
tangent has slope $-3/4$.}
\label{ch13:fig-circle}
\end{figure}

One option is brute force: solve for $y$. Above the $x$-axis the upper half of the circle is the
genuine function $y = \sqrt{25 - x^2}$, and we could differentiate that. It works here, but it is
already a little ugly, and for most equations it is hopeless. Try to solve
\[
x^3 + y^3 = 6xy
\]
for $y$ -- this is the famous \emph{folium of Descartes} -- and you will be staring at the cubic
formula, producing a monstrous expression with three branches. Nobody wants to differentiate that.
We need a method that respects the equation as it stands.

\begin{intuitionbox}
The key mental shift: stop thinking of $y$ as ``the output of a formula'' and start thinking of $y$
as ``whatever value (or values) the equation forces, given $x$.'' On the upper half of the circle,
as $x$ creeps from $3$ toward $4$, the equation \emph{forces} $y$ to creep downward from $4$ toward
$3$. That forced motion of $y$ is exactly what the derivative $\dv{y}{x}$ measures. The equation is
the boss; $y$ obeys.
\end{intuitionbox}

The phrase mathematicians use is that the equation defines $y$ \emph{implicitly} as a function of
$x$. ``Implicit'' means \emph{implied but not spelled out}: the dependence of $y$ on $x$ is real and
present, hiding inside the equation, even though no explicit formula $y=f(x)$ is written down. Our
goal is to extract the slope of that hidden function directly.

\begin{focusbox}
\textbf{Explicit vs. implicit.}
\begin{itemize}
\item \emph{Explicit:} $y = \sqrt{25-x^2}$. The output is isolated; differentiate with the rules you
already know.
\item \emph{Implicit:} $x^2 + y^2 = 25$. The output is entangled with the input. We will
differentiate \emph{both sides of the equation}, treating $y$ as an unknown function of $x$.
\end{itemize}
The whole trick lives in that last phrase.
\end{focusbox}

\section{Implicit Differentiation: the Chain Rule in Disguise}
\label{ch13:implicit}

Here is the central move. We have an equation in $x$ and $y$ that holds along a curve. Along that
curve $y$ is some function of $x$ -- call it $y(x)$ even if we cannot write its formula. Then
\emph{both sides of the equation are functions of $x$}, and they are equal for every $x$ on the
curve. If two functions of $x$ are equal everywhere, their derivatives are equal too. So we may
apply $\dfrac{\dd}{\dd x}$ to both sides and keep a true equation. The only subtlety is that wherever
$y$ appears, it is secretly $y(x)$, so differentiating it requires the chain rule.

Let us do the circle.

\begin{example}[The circle, by implicit differentiation]
\label{ch13:ex-circle}
Differentiate both sides of $x^2 + y^2 = 25$ with respect to $x$:
\[
\ddx{x^2} + \frac{\dd}{\dd x}\!\left(y^2\right) = \frac{\dd}{\dd x}(25).
\]
The first term is $2x$. The right side is $0$, since $25$ is constant. The middle term is the
delicate one. Here $y$ is a function of $x$, so $y^2$ is a \emph{composition}: square-of, applied to
$y(x)$. The chain rule says differentiate the outside ($u^2 \mapsto 2u$) and multiply by the
derivative of the inside ($y$ with respect to $x$):
\[
\frac{\dd}{\dd x}\!\left(y^2\right) = 2y\cdot\dv{y}{x}.
\]
Putting it together,
\[
2x + 2y\,\dv{y}{x} = 0 \quad\Longrightarrow\quad \dv{y}{x} = -\frac{x}{y}.
\]
Now read off the slope at $(3,4)$: $\dv{y}{x} = -\tfrac{3}{4}$. Exactly the geometric prediction.
And the answer at $(-3,-4)$? Also $-\tfrac{3}{4}$ -- the formula handles every branch of the curve at
once, with no case-splitting.
\end{example}

Notice what just happened. We never solved for $y$. The derivative came out in terms of \emph{both}
$x$ and $y$, which is perfectly fine: to use it you supply a point $(x,y)$ that lies on the curve,
and out comes the slope there. An implicit derivative is a slope \emph{field along the curve},
expressed using the coordinates of the point.

\begin{whybox}
Why is differentiating both sides legitimate, and why does $\dv{y}{x}$ pop out? Two facts do all the
work.

\emph{First}, an equation like $x^2+y^2=25$ asserts that two functions of $x$ -- the left side and
the right side, once we substitute $y=y(x)$ -- are equal at every point of the curve. Equal
functions have equal derivatives. (This is the same principle that lets you differentiate both sides
of $\sin^2 x + \cos^2 x = 1$.) So $\dfrac{\dd}{\dd x}(\text{LHS}) = \dfrac{\dd}{\dd x}(\text{RHS})$ is
a valid new equation.

\emph{Second}, every term containing $y$ is really a composition of some operation with the inner
function $y(x)$. The chain rule is precisely the tool for differentiating compositions. So each
$y$-term contributes its outer derivative times the inner derivative $\dv{y}{x}$. Implicit
differentiation is nothing more than ``differentiate the equation, and remember the chain rule
whenever you touch $y$.'' That is why $\dv{y}{x}$ always appears \emph{linearly} (to the first
power, never squared), which is what lets us solve for it by simple algebra at the end.
\end{whybox}

\begin{happeningbox}
When you write $\dfrac{\dd}{\dd x}\!\left(y^2\right) = 2y\,\dv{y}{x}$, what is really happening is a
two-step relay of sensitivities. Nudge $x$ by a tiny $\dd x$. The curve forces $y$ to move by
$\dd y = \dv{y}{x}\,\dd x$. That motion of $y$ then changes $y^2$ by (rate of $y^2$ with respect to
$y$) times $\dd y$, i.e.\ $2y\,\dd y$. Chaining the two responses gives $2y\,\dv{y}{x}\,\dd x$. The
factor $\dv{y}{x}$ is the ``exchange rate'' converting a change in $x$ into the change in $y$ that
the constraint demands. Every appearance of $\dv{y}{x}$ in implicit differentiation is one of these
exchange-rate factors, inserted exactly where the chain of dependence passes through $y$.
\end{happeningbox}

\subsection{A clean recipe}

\begin{focusbox}
\textbf{How to differentiate implicitly.}
\begin{enumerate}
\item Apply $\dfrac{\dd}{\dd x}$ to \emph{both} sides of the equation.
\item Differentiate terms in $x$ normally. For any term involving $y$, use the chain rule: after
differentiating with respect to $y$, multiply by $\dv{y}{x}$. (Products like $xy$ also need the
product rule.)
\item Collect every term containing $\dv{y}{x}$ on one side, everything else on the other.
\item Factor out $\dv{y}{x}$ and divide. The answer will usually involve both $x$ and $y$.
\end{enumerate}
\end{focusbox}

\begin{trapbox}
The single most common mistake is forgetting the $\dv{y}{x}$ factor and writing
$\dfrac{\dd}{\dd x}(y^2) = 2y$. That is wrong: $2y$ is the derivative of $y^2$ \emph{with respect to
$y$}, not with respect to $x$. The whole point is that $y$ depends on $x$, so the chain rule supplies
the extra factor $\dv{y}{x}$. A good habit: every time your pen leaves a $y$-term, immediately tack
on $\dv{y}{x}$ before you forget. By contrast, $\dfrac{\dd}{\dd x}(x^2)=2x$ needs \emph{no} extra
factor, because $\dv{x}{x}=1$.
\end{trapbox}

\begin{example}[A term that needs the product rule]
\label{ch13:ex-product}
Find $\dv{y}{x}$ for $x^2 y + y^3 = 10$.

Differentiate term by term. The term $x^2 y$ is a product of $x^2$ and $y$, so use the product rule,
and remember the chain-rule factor on $y$:
\[
\frac{\dd}{\dd x}\!\left(x^2 y\right) = 2x\cdot y + x^2\cdot \dv{y}{x}.
\]
For $y^3$ the chain rule gives $3y^2\,\dv{y}{x}$, and the right side is $0$. So
\[
2xy + x^2\,\dv{y}{x} + 3y^2\,\dv{y}{x} = 0.
\]
Collect the $\dv{y}{x}$ terms:
\[
\left(x^2 + 3y^2\right)\dv{y}{x} = -2xy
\quad\Longrightarrow\quad
\dv{y}{x} = \frac{-2xy}{x^2 + 3y^2}.
\]
\end{example}

\begin{labbox}
Try one with your own pen before reading on: find $\dv{y}{x}$ for $\sin(xy) = x$. You will need the
chain rule on $\sin$, then the product rule on $xy$. (Solution: differentiating,
$\cos(xy)\cdot\big(y + x\,\dv{y}{x}\big) = 1$, so
$\dv{y}{x} = \dfrac{1 - y\cos(xy)}{x\cos(xy)}$.) The structure is always the same: $\dv{y}{x}$ shows
up linearly, and you solve for it.
\end{labbox}

\subsection{The folium of Descartes}

Now the curve we could not solve. The folium $x^3 + y^3 = 6xy$ is a looping cubic that crosses
itself at the origin -- a shape that genuinely has no single-formula description, yet implicit
differentiation handles it without breaking a sweat.

\begin{example}[Slope on the folium]
\label{ch13:ex-folium}
Differentiate $x^3 + y^3 = 6xy$ with respect to $x$. The right side $6xy$ is a product, so it needs
the product rule (and the chain factor on $y$):
\[
3x^2 + 3y^2\,\dv{y}{x} = 6\!\left(y + x\,\dv{y}{x}\right).
\]
Expand and gather $\dv{y}{x}$ terms:
\[
3y^2\,\dv{y}{x} - 6x\,\dv{y}{x} = 6y - 3x^2,
\qquad
\left(3y^2 - 6x\right)\dv{y}{x} = 6y - 3x^2.
\]
Divide:
\[
\dv{y}{x} = \frac{6y - 3x^2}{3y^2 - 6x} = \frac{2y - x^2}{y^2 - 2x}.
\]
The point $(3,3)$ lies on the curve (check: $27 + 27 = 54 = 6\cdot 9$). There the slope is
\[
\dv{y}{x}\bigg|_{(3,3)} = \frac{6 - 9}{9 - 6} = \frac{-3}{3} = -1.
\]
The tangent there is the line $y - 3 = -1(x-3)$, i.e.\ $y = -x + 6$.
\end{example}

\begin{figure}[H]
\centering
\begin{tikzpicture}
\begin{axis}[
  width=11cm, height=9cm,
  axis lines=middle,
  xlabel={$x$}, ylabel={$y$},
  xmin=-3.2, xmax=4.5, ymin=-3.2, ymax=4.5,
  xtick={-2,2,4}, ytick={-2,2,4},
  samples=200,
]
% folium parametrization: x=6t/(1+t^3), y=6t^2/(1+t^3)
\addplot[curveblue, thick, domain=-0.85:-0.1] ({6*x/(1+x^3)}, {6*x^2/(1+x^3)});
\addplot[curveblue, thick, domain=-0.1:20, samples=200] ({6*x/(1+x^3)}, {6*x^2/(1+x^3)});
\addplot[curveblue, thick, domain=-20:-1.18, samples=200] ({6*x/(1+x^3)}, {6*x^2/(1+x^3)});
% point (3,3) and tangent y=-x+6
\addplot[only marks, mark=*, tangentred] coordinates {(3,3)};
\addplot[tangentred, thick, domain=1.2:4.5] {-x+6};
\node at (axis cs: 3.5,3.3) {$(3,3)$};
\node[tangentred] at (axis cs: 1.1,4.2) {tangent};
\node[curveblue] at (axis cs: 2.4,1.0) {folium};
\end{axis}
\end{tikzpicture}
\caption{The folium of Descartes $x^3+y^3=6xy$. There is no tidy $y=f(x)$ for this looping curve, but
implicit differentiation gives the tangent slope at $(3,3)$ as $-1$.}
\label{ch13:fig-folium}
\end{figure}

\begin{historybox}
Ren\'e Descartes proposed this curve in 1638 in a challenge to Pierre de Fermat, almost as a taunt:
find the tangent lines to \emph{this}. At the time there was no general derivative; Fermat had his own
ingenious ``adequality'' method for tangents and maxima, a precursor of the limit idea in which he
nudged the variable by a small amount, simplified, and then discarded the leftover small term.
Descartes had his own algebraic tangent method using circles that touch the curve. Both could, with
effort, find tangents to the folium. What neither had was the unifying machinery we now use in three
lines. The chain rule -- and with it implicit differentiation -- arrived with Leibniz in the 1670s
and 1680s, whose $\dd y$, $\dd x$ notation was built precisely so that relationships between
differentials could be manipulated algebraically. The very fact that $\dv{y}{x}$ falls out so
cleanly here is a monument to how much labor good notation can save.
\end{historybox}

\section{Tangent Lines and Higher Derivatives}
\label{ch13:tangents}

Once you can find $\dv{y}{x}$ at a point on an implicit curve, tangent lines cost nothing extra: the
tangent at $(x_0,y_0)$ is
\[
y - y_0 = m\,(x - x_0), \qquad m = \dv{y}{x}\bigg|_{(x_0,y_0)}.
\]
The only thing to watch is that the point must actually lie on the curve, and that you substitute
\emph{both} coordinates into the slope formula.

\begin{example}[Tangent and a vertical tangent on an ellipse]
\label{ch13:ex-ellipse}
Consider the ellipse $\dfrac{x^2}{9} + \dfrac{y^2}{4} = 1$. Differentiate implicitly:
\[
\frac{2x}{9} + \frac{2y}{4}\,\dv{y}{x} = 0
\quad\Longrightarrow\quad
\dv{y}{x} = -\frac{4x}{9y}.
\]
At the point $\big(\tfrac{3}{\sqrt2}, \sqrt2\big)$ (which lies on the ellipse: $\tfrac{9/2}{9} +
\tfrac{2}{4} = \tfrac12 + \tfrac12 = 1$), the slope is
\[
m = -\frac{4\cdot \tfrac{3}{\sqrt2}}{9\cdot \sqrt2} = -\frac{12/\sqrt2}{9\sqrt2}
= -\frac{12}{9\cdot 2} = -\frac{2}{3}.
\]
Notice also \emph{where the tangent is vertical}: the slope formula $-\tfrac{4x}{9y}$ blows up when
$y=0$, i.e.\ at $(\pm 3, 0)$, the ends of the major axis. There the tangent line is vertical -- a
slope our $\dv{y}{x}$ honestly reports as ``undefined,'' which is the correct geometric answer.
Implicit differentiation does not hide these features; it flags them
(Figure~\ref{ch13:fig-ellipse}).
\end{example}

\begin{figure}[H]
\centering
\begin{tikzpicture}
\begin{axis}[
  width=12cm, height=8cm,
  axis equal image,
  axis lines=middle,
  xlabel={$x$}, ylabel={$y$},
  xmin=-4.3, xmax=4.3, ymin=-3.2, ymax=3.2,
  xtick={-3,3}, ytick={-2,2},
  samples=120,
]
% ellipse x^2/9 + y^2/4 = 1, parametric
\addplot[curveblue, thick, domain=0:360] ({3*cos(x)}, {2*sin(x)});
% slanted tangent at (3/sqrt2, sqrt2): slope -2/3
\addplot[tangentred, thick, domain=-0.3:4.0] {1.41421 - (2/3)*(x-2.12132)};
\addplot[only marks, mark=*, tangentred] coordinates {(2.12132,1.41421)};
% vertical tangents at (+/-3, 0)
\addplot[structurepurple, thick, dashed] coordinates {(3,-2.6) (3,2.6)};
\addplot[structurepurple, thick, dashed] coordinates {(-3,-2.6) (-3,2.6)};
\addplot[only marks, mark=*, structurepurple] coordinates {(3,0) (-3,0)};
\node[tangentred] at (axis cs: 3.3,2.4) {slope $-\tfrac23$};
\node[structurepurple] at (axis cs: 3.05,-2.95) {vertical};
\node[structurepurple] at (axis cs: -3.0,-2.95) {vertical};
\end{axis}
\end{tikzpicture}
\caption{The ellipse $\tfrac{x^2}{9}+\tfrac{y^2}{4}=1$. At a typical point the implicit slope
$-\tfrac{4x}{9y}$ gives an ordinary slanted tangent (here $-\tfrac23$). At the ends of the major
axis $(\pm 3,0)$, where $y=0$, the formula blows up and the tangent is vertical -- exactly the
feature that ``solve for $y$'' handles badly.}
\label{ch13:fig-ellipse}
\end{figure}

\begin{connectionbox}
Vertical tangents are where the implicit picture genuinely beats ``solve for $y$.'' If you split the
ellipse into $y = \pm\tfrac{2}{3}\sqrt{9-x^2}$, the derivative of each branch becomes infinite and
ill-behaved exactly at $x=\pm 3$, the seam where the two branches meet. The implicit form
$\dv{y}{x} = -\tfrac{4x}{9y}$ treats the whole curve as one object and simply says: slope undefined
when $y=0$. Later, when we study curves in the plane more generally, this is the seed of the idea
that a smooth curve can have a well-defined \emph{tangent direction} even where it has no
well-defined \emph{function slope}.
\end{connectionbox}

\subsection{Second derivatives implicitly}

You can differentiate a second time to get $\dfrac{\dd^2 y}{\dd x^2}$, useful for concavity. The
method is the same -- differentiate again, remembering the chain rule -- but now you also substitute
the expression for $\dv{y}{x}$ you already found.

\begin{example}[Second derivative on the circle]
\label{ch13:ex-second}
For $x^2+y^2=25$ we found $\dv{y}{x} = -\dfrac{x}{y}$. Differentiate this quotient with respect to
$x$, using the quotient rule and remembering that $y$ depends on $x$:
\[
\frac{\dd^2 y}{\dd x^2}
= -\,\frac{\dfrac{\dd}{\dd x}(x)\cdot y - x\cdot\dfrac{\dd}{\dd x}(y)}{y^2}
= -\,\frac{y - x\,\dv{y}{x}}{y^2}.
\]
Now substitute $\dv{y}{x} = -\tfrac{x}{y}$:
\[
\frac{\dd^2 y}{\dd x^2}
= -\,\frac{y - x\left(-\tfrac{x}{y}\right)}{y^2}
= -\,\frac{y + \tfrac{x^2}{y}}{y^2}
= -\,\frac{y^2 + x^2}{y^3}.
\]
But on the circle $x^2 + y^2 = 25$, so this simplifies beautifully to
\[
\frac{\dd^2 y}{\dd x^2} = -\frac{25}{y^3}.
\]
On the upper half ($y>0$) this is negative -- the curve is concave down, as a dome should be. On the
lower half ($y<0$) it is positive, concave up. The sign of the second derivative reads off the shape
of each half exactly as it should.
\end{example}

\begin{walkbox}
Take this on a walk. The relation $x^2+y^2=25$ is symmetric in $x$ and $y$ -- swapping them leaves
the circle unchanged. We found $\dv{y}{x} = -\tfrac{x}{y}$. What should $\dv{x}{y}$ be, and how is it
related? Convince yourself geometrically that $\dv{x}{y} = -\tfrac{y}{x}$ and that
$\dv{y}{x}\cdot\dv{x}{y} = 1$, then explain in words why ``run over rise'' is the reciprocal of
``rise over run.'' This is the first whisper of the inverse-function relationship between derivatives.
\end{walkbox}

\section{Related Rates: When One Rate Drives Another}
\label{ch13:related-rates}

We now turn the same idea toward \emph{time}. In a related-rates problem, several quantities change
as time passes, and they are tied together by a geometric or physical relationship. Because the
quantities are bound by an equation, their \emph{rates of change} are bound too -- and if we know
some of the rates we can solve for the others. The mathematics is implicit differentiation again,
but now the hidden variable everything depends on is $t$, time.

\begin{intuitionbox}
Picture a constraint as a length of inextensible string connecting two beads. If you know how fast
one bead moves, the string \emph{forces} a definite speed on the other; you do not get to choose it.
A related-rates problem hands you the constraint (the string) and one or more known rates, and asks
for an unknown rate. The constraint, differentiated in time, is the rule that converts the rates you
know into the rate you want.
\end{intuitionbox}

The crucial structural fact: every quantity in the problem is a function of time. So when we
differentiate the constraint equation with respect to $t$, \emph{every variable} gets a chain-rule
factor -- its own rate $\dfrac{\dd(\cdot)}{\dd t}$. A length $x$ contributes $\dv{x}{t}$, an area $A$
contributes $\dv{A}{t}$, an angle $\theta$ contributes $\dv{\theta}{t}$. There is no privileged
``independent variable'' among $x,y,\theta$; they are all dependent on the one true independent
variable, $t$.

\begin{happeningbox}
Here is the real mechanism, stripped bare. A constraint $F(x,y) = c$ holds at every instant. As time
ticks, the point $(x(t), y(t))$ slides along the curve $F=c$; it is not free to wander off. The
condition for staying on the curve is that the \emph{rate} at which $F$ changes must be zero --
$\dv{F}{t} = 0$ -- because $F$ is pinned at the constant $c$. Writing out $\dv{F}{t}=0$ by the chain
rule produces a single linear equation relating $\dv{x}{t}$ and $\dv{y}{t}$. That one equation
\emph{is} the related-rates relationship. ``One rate drives another'' is the statement that this
equation has only one degree of freedom: fix one rate and the other is determined. Related rates is
the time-derivative of a constraint, and nothing more.
\end{happeningbox}

\subsection{A modeling method that always works}

Related-rates problems feel hard because they come wrapped in a story. The cure is a fixed procedure
that turns the story into algebra in the same way every time.

\begin{focusbox}
\textbf{The five-step method.}
\begin{enumerate}
\item \textbf{Draw and name.} Sketch the situation. Assign a \emph{variable} to every quantity that
changes. Use a letter, not a number, even for quantities you happen to know right now -- their
numerical value is only true at one instant, and plugging it in too early destroys the relationship.
\item \textbf{List knowns and the unknown rate.} Translate every ``how fast'' into a derivative with
respect to $t$. Watch the signs: a shrinking quantity has a negative rate.
\item \textbf{Find the constraint.} Write a single equation connecting your variables, valid at
\emph{all} times -- usually a geometric fact (Pythagoras, similar triangles, area or volume formula,
trig ratio).
\item \textbf{Differentiate with respect to $t$.} Every variable earns its own $\dfrac{\dd}{\dd t}$
factor (the chain rule, as always). This gives an equation relating the rates.
\item \textbf{Substitute and solve.} \emph{Now}, and only now, plug in the instantaneous numbers and
solve for the unknown rate. Interpret the sign and units.
\end{enumerate}
\end{focusbox}

\begin{trapbox}
\textbf{Do not substitute numbers before differentiating.} If you replace a changing length by its
current value \emph{before} taking $\dfrac{\dd}{\dd t}$, you have turned a moving quantity into a
constant, and its derivative becomes a wrong $0$. The numerical values describe a single frozen
instant; the constraint equation and its differentiation describe \emph{all} instants. Differentiate
first (in symbols), substitute last (with numbers). This one discipline prevents the majority of
related-rates errors.
\end{trapbox}

\subsection{The sliding ladder}

\begin{example}[The ladder against the wall]
\label{ch13:ex-ladder}
A $10$-foot ladder leans against a vertical wall. The bottom is pulled away from the wall at
$2$\,ft/s. How fast is the top sliding down the wall when the bottom is $6$\,ft from the wall?

\textbf{Draw and name.} Let $x$ be the distance from the wall to the foot of the ladder, and $y$ the
height of the top on the wall. The ladder's length is the fixed constraint.

\begin{figure}[H]
\centering
\begin{tikzpicture}[scale=0.9]
% wall and ground
\draw[gray!60, thick] (0,0) -- (0,5.2) node[above] {wall};
\draw[gray!60, thick] (0,0) -- (5.5,0) node[right] {ground};
% ladder
\draw[curveblue, very thick] (3,0) -- (0,4) ;
\node[curveblue] at (1.9,2.4) {ladder, $10$};
% labels
\draw[<->, structurepurple] (0,-0.35) -- (3,-0.35) node[midway, below] {$x$};
\draw[<->, structurepurple] (-0.4,0) -- (-0.4,4) node[midway, left] {$y$};
% foot moving right
\draw[->, tangentred, thick] (3.1,0.0) -- (3.9,0.0) node[right] {$\dv{x}{t}=2$};
% top moving down
\draw[->, tangentred, thick] (0.0,3.9) -- (0.0,3.1) node[left, yshift=-2pt] {$\dv{y}{t}=?$};
\draw[fill=black] (3,0) circle (1.5pt);
\draw[fill=black] (0,4) circle (1.5pt);
\end{tikzpicture}
\caption{The sliding ladder. The fixed length $10$ ties $x$ and $y$ together by
$x^2+y^2=100$; pulling the foot out forces the top down.}
\label{ch13:fig-ladder}
\end{figure}

\textbf{Knowns and unknown.} $\dv{x}{t} = +2$\,ft/s (the foot moves away). We want $\dv{y}{t}$ at the
instant $x=6$.

\textbf{Constraint.} The ladder is rigid, so by the Pythagorean theorem, \emph{at every instant},
\[
x^2 + y^2 = 100.
\]
\textbf{Differentiate in $t$.} Each variable gets its rate:
\[
2x\,\dv{x}{t} + 2y\,\dv{y}{t} = 0
\quad\Longrightarrow\quad
\dv{y}{t} = -\frac{x}{y}\dv{x}{t}.
\]
\textbf{Substitute.} When $x=6$, the constraint gives $y = \sqrt{100-36} = 8$. So
\[
\dv{y}{t} = -\frac{6}{8}\cdot 2 = -\frac{3}{2} = -1.5\ \text{ft/s}.
\]
The top slides \emph{down} at $1.5$\,ft/s. The negative sign is the model telling us the height is
decreasing, exactly as it must.
\end{example}

\begin{aestheticbox}
There is a startling rider to the ladder problem. Keep the foot moving at a steady $2$\,ft/s and ask
what happens as the ladder nears flat against the ground, $y\to 0$. Then $\dv{y}{t} = -\tfrac{x}{y}
\cdot 2$ blows up to $-\infty$: the top appears to slam down infinitely fast. Of course no real
ladder does this -- near the end the top lifts off the wall and the model's assumptions fail. But the
mathematics is honest about its own breakdown: the constraint $x^2+y^2=100$ forces an infinite rate,
which is the model's way of shouting that something physical (the ladder leaving the wall) has
intervened. Good models tell you when to stop trusting them.
\end{aestheticbox}

\subsection{The inflating balloon}

\begin{example}[The spherical balloon]
\label{ch13:ex-balloon}
Air is pumped into a spherical balloon at $100$\,cm$^3$/s. How fast is the radius growing when the
radius is $5$\,cm? How fast is it growing when the radius is $10$\,cm?

\textbf{Name.} Let $V$ be the volume and $r$ the radius, both functions of $t$.

\textbf{Knowns and unknown.} $\dv{V}{t} = 100$\,cm$^3$/s. Want $\dv{r}{t}$.

\textbf{Constraint.} The geometry of a sphere:
\[
V = \frac{4}{3}\pi r^3.
\]
\textbf{Differentiate in $t$} (chain rule on $r^3$):
\[
\dv{V}{t} = 4\pi r^2\,\dv{r}{t}
\quad\Longrightarrow\quad
\dv{r}{t} = \frac{1}{4\pi r^2}\dv{V}{t}.
\]
\textbf{Substitute.} At $r=5$:
\[
\dv{r}{t} = \frac{100}{4\pi(25)} = \frac{100}{100\pi} = \frac{1}{\pi}\approx 0.318\ \text{cm/s}.
\]
At $r=10$:
\[
\dv{r}{t} = \frac{100}{4\pi(100)} = \frac{1}{4\pi}\approx 0.0796\ \text{cm/s}.
\]
\end{example}

\begin{intuitionbox}
The factor $4\pi r^2$ in $\dv{V}{t} = 4\pi r^2\,\dv{r}{t}$ is not a coincidence -- it is the
\emph{surface area} of the sphere. The air you pump in spreads itself over the skin of the balloon as
a thin shell; a thicker shell over a bigger skin means the same volume of air buys less radius. That
is precisely why the radius slows as the balloon grows: at $r=10$ the surface is four times that at
$r=5$, so the same $100$\,cm$^3$/s thickens it only a quarter as fast. The derivative $\dv{V}{r} =
4\pi r^2 = (\text{surface area})$ is one of the most satisfying ``of course!'' facts in calculus.
\end{intuitionbox}

\subsection{The walking shadow}

\begin{example}[The lengthening shadow]
\label{ch13:ex-shadow}
A street lamp is $15$\,ft tall. A person $6$\,ft tall walks away from the base of the lamp at
$5$\,ft/s. (a) How fast is the tip of the shadow moving? (b) How fast is the shadow's length growing?

\textbf{Draw and name.} Let $x$ be the person's distance from the lamp post, and $s$ the length of
the shadow beyond the person. The tip of the shadow is then at distance $x+s$ from the post.

\begin{figure}[H]
\centering
\begin{tikzpicture}[scale=0.85]
% ground
\draw[gray!60, thick] (-0.3,0) -- (10.5,0);
% lamp post
\draw[curveblue, very thick] (0,0) -- (0,5) ;
\node[curveblue, above] at (0,5) {$15$};
\draw[fill=yellow!70!orange] (0,5) circle (3pt);
% person
\draw[structurepurple, very thick] (5,0) -- (5,2);
\node[structurepurple, right] at (5,1) {$6$};
% light rays
\draw[orange!70!black] (0,5) -- (8,0);
\draw[orange!70!black, dashed] (0,5) -- (5,2);
% labels
\draw[<->] (0,-0.4) -- (5,-0.4) node[midway,below] {$x$};
\draw[<->] (5,-0.4) -- (8,-0.4) node[midway,below] {$s$};
\draw[->, tangentred, thick] (5.1,2.0) -- (6.0,2.0) node[right] {$5$ ft/s};
\draw[fill=black] (8,0) circle (1.5pt);
\node[below right] at (8,0) {tip};
\end{tikzpicture}
\caption{The shadow problem. Similar triangles -- the big one from lamp top to shadow tip, the small
one from the person's head to the tip -- give the constraint.}
\label{ch13:fig-shadow}
\end{figure}

\textbf{Constraint.} The lamp, the person, and the two shadow tips form \emph{similar triangles}: the
large triangle (height $15$, base $x+s$) and the small triangle (height $6$, base $s$) share the same
shape. Equal ratios of height to base:
\[
\frac{15}{x+s} = \frac{6}{s}.
\]
Cross-multiply: $15 s = 6(x+s) = 6x + 6s$, so $9s = 6x$, giving the clean linear relation
\[
s = \frac{2}{3}x.
\]
\textbf{Differentiate in $t$.}
\[
\dv{s}{t} = \frac{2}{3}\dv{x}{t} = \frac{2}{3}(5) = \frac{10}{3}\ \text{ft/s}.
\]
That answers (b): the shadow lengthens at $\tfrac{10}{3}\approx 3.33$\,ft/s.

For (a), the tip is at position $x+s$, so its speed is
\[
\dv{}{t}(x+s) = \dv{x}{t} + \dv{s}{t} = 5 + \frac{10}{3} = \frac{25}{3}\approx 8.33\ \text{ft/s}.
\]
\end{example}

\begin{trapbox}
Part (a) and part (b) ask different questions, and confusing them is the classic shadow trap. The
\emph{length of the shadow} ($s$) and the \emph{position of its tip} ($x+s$) are not the same
quantity; the tip moves because \emph{both} the person and the shadow's far end advance. Always
decide precisely which point or length the problem is asking about before differentiating. Notice
too that neither rate depends on $x$ -- the answers are the same whether the person is near or far.
That is special to the straight-line, point-source geometry here, and it is a nice surprise worth
pausing on.
\end{trapbox}

\subsection{An angle that changes: the rising balloon}

\begin{example}[Tracking with a camera]
\label{ch13:ex-angle}
A weather balloon rises straight up at $8$\,m/s. A camera on the ground sits $50$\,m from the launch
point and tilts to keep the balloon in frame. How fast is the camera's angle of elevation increasing
when the balloon is $50$\,m high?

\textbf{Name.} Let $h$ be the balloon's height and $\theta$ the camera's angle of elevation above the
horizontal. The horizontal distance $50$\,m is constant.

\begin{figure}[H]
\centering
\begin{tikzpicture}[scale=0.9]
% ground
\draw[gray!60, thick] (-0.3,0) -- (6.0,0);
% camera at origin, launch point at (5,0)
\draw[fill=black] (0,0) circle (1.6pt);
\node[below left] at (0,0) {camera};
% vertical path of balloon at x=5
\draw[structurepurple, thick] (5,0) -- (5,3.2);
\draw[fill=yellow!70!orange] (5,3.2) circle (3pt);
\node[structurepurple, right] at (5,1.6) {$h$};
% line of sight
\draw[curveblue, thick] (0,0) -- (5,3.2);
% horizontal base label
\draw[<->] (0,-0.45) -- (5,-0.45) node[midway,below] {$50$};
% angle
\draw[tangentred] (1.1,0) arc (0:32.6:1.1);
\node[tangentred] at (1.5,0.35) {$\theta$};
% rising arrow
\draw[->, tangentred, thick] (5.0,3.3) -- (5.0,3.9) node[right] {$8$ m/s};
\end{tikzpicture}
\caption{Tracking the rising balloon. The camera sits a fixed $50$\,m from the launch point; the
line of sight makes angle $\theta$ with the ground, and $\tan\theta = h/50$. As $h$ grows, $\theta$
opens.}
\label{ch13:fig-camera}
\end{figure}

\textbf{Constraint.} From the right triangle (Figure~\ref{ch13:fig-camera}),
\[
\tan\theta = \frac{h}{50}.
\]
\textbf{Differentiate in $t$.} The left side needs the chain rule on $\tan$:
\[
\sec^2\theta\,\dv{\theta}{t} = \frac{1}{50}\dv{h}{t}.
\]
\textbf{Substitute.} When $h=50$, the triangle is a $50$-$50$ right triangle, so $\theta = 45^\circ$
and the hypotenuse is $50\sqrt2$. Then $\sec\theta = \dfrac{50\sqrt2}{50} = \sqrt2$, so
$\sec^2\theta = 2$. With $\dv{h}{t}=8$,
\[
2\,\dv{\theta}{t} = \frac{8}{50} = \frac{4}{25}
\quad\Longrightarrow\quad
\dv{\theta}{t} = \frac{2}{25} = 0.08\ \text{rad/s}.
\]
The angle increases at $0.08$\,rad/s, about $4.6$ degrees per second.
\end{example}

\begin{trapbox}
The factor $\sec^2\theta$ is easy to drop. Because $\theta$ changes with time, $\tan\theta$ is a
composition (tangent-of, applied to $\theta(t)$), and the chain rule supplies
$\dfrac{\dd}{\dd t}\tan\theta = \sec^2\theta\,\dv{\theta}{t}$. Forgetting it is the trig version of
the $\dv{y}{x}$-dropping mistake from earlier. Also: this calculation only works because $\theta$ is
measured in \emph{radians} -- the derivative formula $\dfrac{\dd}{\dd\theta}\tan\theta = \sec^2\theta$
is a radian fact. Always work related-rates angles in radians.
\end{trapbox}

\subsection{Two moving pieces: ships on the sea}

So far one rate drove one other. The method is just as happy with several changing lengths at once.

\begin{example}[Two ships]
\label{ch13:ex-ships}
At noon, ship $A$ is $60$\,km due east of ship $B$. Ship $A$ sails west at $20$\,km/h; ship $B$ sails
north at $15$\,km/h. How fast is the distance between them changing at $1$\,p.m.?

\textbf{Name.} Put $B$'s noon position at the origin. Let $p$ be the east--west separation between
the ships and $q$ the north--south separation, and $D$ the straight-line distance between them.

\textbf{Set up.} Ship $A$ closes the east--west gap at $20$\,km/h, so $p = 60 - 20t$ and
$\dv{p}{t} = -20$. Ship $B$ opens a north--south gap at $15$\,km/h, so $q = 15t$ and $\dv{q}{t} =
15$. (Here $t$ is hours after noon.)

\textbf{Constraint.} The two ships and the right angle between their motions give
\[
D^2 = p^2 + q^2.
\]
\textbf{Differentiate in $t$.}
\[
2D\,\dv{D}{t} = 2p\,\dv{p}{t} + 2q\,\dv{q}{t}
\quad\Longrightarrow\quad
\dv{D}{t} = \frac{p\,\dv{p}{t} + q\,\dv{q}{t}}{D}.
\]
\textbf{Substitute at $t=1$.} Then $p = 60-20 = 40$, $q = 15$, and $D = \sqrt{40^2 + 15^2} =
\sqrt{1600 + 225} = \sqrt{1825} \approx 42.72$\,km. So
\[
\dv{D}{t} = \frac{40(-20) + 15(15)}{\sqrt{1825}} = \frac{-800 + 225}{42.72}
= \frac{-575}{42.72}\approx -13.5\ \text{km/h}.
\]
The ships are still \emph{approaching} at about $13.5$\,km/h: the negative sign says the distance is
shrinking, because $A$'s rapid westward closing outweighs $B$'s slower northward drift -- for now.
\end{example}

\begin{walkbox}
The ships' distance is decreasing at $1$\,p.m., but $A$ will eventually slide past $B$'s line and the
gap must grow again. So somewhere $\dv{D}{t}$ passes through $0$ -- the moment of closest approach.
Without solving it fully, can you see that the closest approach happens exactly when $p\,\dv{p}{t} +
q\,\dv{q}{t} = 0$? That combination is a dot product in disguise: the rate of $D$ is the component of
the relative velocity along the line joining the ships. Related rates is quietly teaching you vector
geometry.
\end{walkbox}

\section{What Ties It All Together}
\label{ch13:synthesis}

Step back and the two halves of this chapter are one idea seen from two angles.

In \emph{implicit differentiation}, a constraint $F(x,y)=c$ binds $y$ to $x$. Differentiating in $x$
and remembering the chain rule on $y$ yields a relationship between $1$ (the rate of $x$ with respect
to itself) and $\dv{y}{x}$, which we solve for the slope.

In \emph{related rates}, a constraint $F(x,y)=c$ binds $x$ and $y$ as both ride along on time.
Differentiating in $t$ and remembering the chain rule on every variable yields a relationship between
$\dv{x}{t}$ and $\dv{y}{t}$, which we solve for the unknown rate.

\begin{aestheticbox}
Same equation, same chain rule, same algebra at the end -- only the independent variable differs ($x$
in one case, $t$ in the other). In fact you can pass between them: if $x$ and $y$ both depend on $t$,
then by the chain rule
\[
\dv{y}{t} = \dv{y}{x}\cdot\dv{x}{t}.
\]
The implicit slope $\dv{y}{x}$ is exactly the exchange rate that converts the rate of $x$ into the
rate of $y$. Implicit differentiation finds the exchange rate; related rates uses it. The deep theme
of this whole chapter is a single sentence: \emph{when quantities obey an equation, their rates obey
the differentiated equation.} Constraints on values become constraints on rates. That principle will
return, in grander dress, when we reach the gradient and the method of Lagrange multipliers --
where the same ``the constraint pins the rate'' logic governs optimization on curved surfaces.
\end{aestheticbox}

\begin{historybox}
This rate-chaining is the beating heart of Newtonian physics, and it is no accident that Newton's own
calculus was built around it. Newton called the variables \emph{fluents} (flowing quantities) and
their time-rates \emph{fluxions}, writing $\dot{x}$ for what we write $\dv{x}{t}$. His central problem
was exactly the related-rates problem: \emph{given a relation among fluents, find the relation among
their fluxions, and vice versa.} The dot notation survives today in mechanics, where $\dot{x}$ is
velocity and $\ddot{x}$ acceleration. When you differentiate $x^2+y^2=100$ in time to get $x\dot{x} +
y\dot{y} = 0$, you are doing precisely what Newton did in the 1660s -- and the fact that it still
reads cleanly three and a half centuries later is a tribute to how natural the idea is once you see
the constraint as the thing that flows.
\end{historybox}

\section{Exercises}
\label{ch13:exercises}

\begin{exercise}
Find $\dv{y}{x}$ by implicit differentiation for each curve.
\begin{enumerate}[label=(\alph*)]
\item $x^2 + xy + y^2 = 7$.
\item $\sqrt{x} + \sqrt{y} = 4$.
\item $e^{xy} = x + y$.
\item $y\sin x = x\cos y$.
\end{enumerate}
\end{exercise}

\begin{exercise}
The curve $x^2 + y^2 = 25$ passes through $(-4,3)$. Find the equation of the tangent line there, and
confirm with the formula $\dv{y}{x} = -\tfrac{x}{y}$ that it is perpendicular to the radius to that
point.
\end{exercise}

\begin{exercise}
For the astroid $x^{2/3} + y^{2/3} = 4$, find $\dv{y}{x}$. At which points is the tangent line
horizontal, and at which is it vertical? (Hint: horizontal where $\dv{y}{x}=0$, vertical where it is
undefined.)
\end{exercise}

\begin{exercise}
Find the second derivative $\dfrac{\dd^2 y}{\dd x^2}$ for the curve $xy = 1$, first by implicit
differentiation and then by solving $y = 1/x$ explicitly and differentiating twice. Confirm the two
answers agree.
\end{exercise}

\begin{exercise}[Conceptual]
A classmate computes $\dfrac{\dd}{\dd x}(y^3) = 3y^2$ while differentiating an implicit equation.
Explain in one or two sentences exactly what is wrong, and what the correct expression is. Then state
the analogous mistake one might make in a related-rates problem with a term $r^3$ when differentiating
in $t$.
\end{exercise}

\begin{exercise}
Water drains from an inverted right-circular cone (point down) of height $12$\,m and top radius
$3$\,m. The volume of water at depth $h$ is $V = \tfrac{1}{3}\pi r^2 h$, where $r$ is the surface
radius at that depth. (a) Use similar triangles to express $r$ in terms of $h$, and show $V =
\tfrac{\pi}{48}h^3$. (b) If the volume decreases at $2$\,m$^3$/min, how fast is the depth dropping
when $h = 4$\,m?
\end{exercise}

\begin{exercise}
A kite $100$\,ft above the ground moves horizontally at $8$\,ft/s away from the person flying it. The
string is taut and runs in a straight line. How fast is the string being let out when $200$\,ft of
string is out? (Hint: let $x$ be the horizontal distance and $L$ the string length, with $x^2 +
100^2 = L^2$.)
\end{exercise}

\begin{exercise}
Two roads meet at a right angle. Car $A$ is $0.3$\,km from the intersection driving toward it at
$60$\,km/h; car $B$ is $0.4$\,km from the intersection driving toward it at $80$\,km/h. How fast is
the distance between the cars changing at this instant? Is it increasing or decreasing, and what does
the sign mean physically?
\end{exercise}

\begin{exercise}
A trough $10$\,ft long has a cross-section that is an isosceles triangle, $2$\,ft across the top and
$1$\,ft deep. Water is poured in at $3$\,ft$^3$/min. How fast is the water level rising when the water
is $6$\,inches deep? (Set up similar triangles for the water's surface width in terms of depth, then
relate volume to depth.)
\end{exercise}

\begin{exercise}[Proof]
Suppose $x$ and $y$ both depend on $t$ and satisfy the constraint $F(x,y)=c$ for all $t$, with $y$
expressible as a differentiable function of $x$ along the curve. Using only the single-variable chain
rule, prove that
\[
\dv{y}{t} = \dv{y}{x}\cdot\dv{x}{t}
\]
at any point where these derivatives exist. Explain in one sentence how this identity unifies
implicit differentiation and related rates.
\end{exercise}

\begin{exercise}[Conceptual]
In the sliding-ladder example we found $\dv{y}{t}\to -\infty$ as the top approaches the ground.
Explain physically why a real ladder does not behave this way, and identify which modeling assumption
($x^2 + y^2 = 100$ holding for all time, the foot moving at constant $2$\,ft/s, the top staying in
contact with the wall) is the first to fail. Argue that the infinite rate is the model honestly
reporting its own breakdown rather than a real physical infinity.
\end{exercise}
